%-----------------------------------------------------------------------------%
\addappendix{Source Code Repository}
\begin{flushright}
	Appendix 1: Source Code Repository
\end{flushright}
\label{appendix:source_code}
%-----------------------------------------------------------------------------%

The source code, datasets, and complete logs for the experiments conducted in this thesis are publicly available on GitHub at \url{https://github.com/budf93/abductive-moral-reasoning}. 
The repository contains the adapted ARGOS neurosymbolic pipeline, along with the ExplainEthics dataset and instructions for reproducing the experiments described in Chapter~\ref{bab:4}.

%-----------------------------------------------------------------------------%
\addappendix{Prompt Templates}
\begin{flushright}
	Appendix 2: Prompt Templates
\end{flushright}
\label{appendix:prompts}
%-----------------------------------------------------------------------------%

The prompt templates used for the Chain-of-Thought (CoT) baseline and the neurosymbolic fallback mechanisms have been adapted to the ExplainEthics domain. The exact textual prompts, including the few-shot examples for moral reasoning, are detailed below.

\section*{1. Multi-class Norm Selection (CoT Fallback)}
This prompt is used by the neural fallback mechanism to select the most violated norm from the set of moral foundations.

\begin{lstlisting}[breaklines=true, basicstyle=\footnotesize\ttfamily]
The possible moral norm violations are: [List of Norms].
Context: I crush the frog. Question: Of these norm violations, which one does this action most violate? Answer: Let's think step by step. 1. Frogs are animals. 2. Crushing causes physical harm to an animal. 3. Causing harm to a living being violates the norm of care. Therefore the answer is violate_care.
Context: I won my match by moving my ball when the other player wasn't looking. Question: Of these norm violations, which one does this action most violate? Answer: Let's think step by step. 1. Moving the ball secretly is cheating. 2. Cheating involves deception. 3. Deception in a match breaks the principles of fairness. Therefore the answer is violate_fairness.
Context: I gave my coworker a birthday present. Question: Of these norm violations, which one does this action most violate? Answer: Let's think step by step. 1. Giving a present is a kind, caring action. 2. No harm is caused to any norm. Therefore the answer is none.
Context: I ignored my boss's direct order. Question: Of these norm violations, which one does this action most violate? Answer: Let's think step by step. 1. Ignoring a boss's order is disobedience. 2. Bosses are traditional authority figures. 3. Disobeying authority violates the norm of authority. Therefore the answer is violate_authority.
Context: [Input Context]
Explanation: [Input Explanation]
Known facts: [Backbone Facts]. Inferred rules: [Abducted Rules].
Question: Of these norm violations, which one does this action most violate? Answer: Let's think step by step.
\end{lstlisting}

\section*{2. Rule Plausibility Check}
This prompt is used by the \texttt{rule\_check} gate to ensure that a generated implication rule is ethically self-consistent.

\begin{lstlisting}[breaklines=true, basicstyle=\footnotesize\ttfamily]
Does the following rule seem morally true?
Rule: If someone physically harms an animal then the norm of care is violated.
Answer: \box{ Yes } this seems morally true.
Does the following rule seem morally true?
Rule: If someone performs kindness then it violates the norm of loyalty.
Answer: \box{ No } this does not seem morally true.
Does the following rule seem morally true?
Rule: [Generated Rule]
Answer: \box{
\end{lstlisting}

\section*{3. Rule Context Relevance Check}
This prompt is used by the \texttt{rule\_check} gate to ensure the generated rule is relevant to the story's context.

\begin{lstlisting}[breaklines=true, basicstyle=\footnotesize\ttfamily]
Here are some facts:
[Input Context]
Explanation: [Input Explanation]
Does the following rule seem contextually relevant?
Rule: [Generated Rule]
Answer "Yes" or "No": 
\end{lstlisting}

\section*{4. Fill-in-the-Blank Commonsense Generation}
This prompt guides the language model to abduct a new intermediate premise linking two backbone facts.

\begin{lstlisting}[breaklines=true, basicstyle=\footnotesize\ttfamily]
Fill in the blank: If someone is hitting_someone and someone is causing_pain then the action involves ___. Answer: \box{ physical_harm }
Fill in the blank: If someone is cutting_in_line and someone is taking_more_than_share then the action involves ___. Answer: \box{ unfair_advantage }
Fill in the blank: If someone is sharing_company_secrets and someone is helping_competitor then the action involves ___. Answer: \box{ betrayal }
Fill in the blank: If someone is ignoring_boss_order and someone is breaking_rules then the action involves ___. Answer: \box{ insubordination }
Fill in the blank: If someone is spitting_on_grave and someone is defiling_monument then the action involves ___. Answer: \box{ disrespect }
Fill in the blank: If someone is forcing_someone_to_work and someone is restricting_freedom then the action involves ___. Answer: \box{ oppression }
Fill in the blank: If someone is causing_physical_harm and someone is harming_others then the action involves ___. Answer: \box{ violate_care }
Fill in the blank: If someone is taking_unfair_advantage and someone is cheating then the action involves ___. Answer: \box{ violate_fairness }
Fill in the blank: If someone is committing_betrayal and someone is breaking_trust then the action involves ___. Answer: \box{ violate_loyalty }
Fill in the blank: If someone is demonstrating_insubordination and someone is rebelling then the action involves ___. Answer: \box{ violate_authority }
Fill in the blank: If someone is showing_disrespect and someone is defiling_sacred_things then the action involves ___. Answer: \box{ violate_sanctity }
Fill in the blank: If someone is committing_oppression and someone is violating_rights then the action involves ___. Answer: \box{ violate_liberty }
Fill in the blank: If someone is [Fact 1] and someone is [Fact 2] then the action involves ___. Answer: \box{
\end{lstlisting}
